% Source-derived chapter 01: Introduction
% Original source: motivic-invariants-birational-maps
\chapter{Introduction}
\label{motivic-invariants-birational-maps-chapter-01}
\source{motivic-invariants-birational-maps}

\begin{remark}[Source section]
\label{motivic-invariants-birational-maps-chapter-01-source}
\source{motivic-invariants-birational-maps}
\source{motivic-invariants-birational-maps}
This chapter reproduces the corresponding section of the retrieved
LaTeX source, including its definitions, statements, examples, and
proofs. It has not yet been translated into Lean.
\notready
\end{remark}

% The source section title was: Introduction
Let $\k$ be a field.
The Cremona groups $\Cr_n(\k) = \Bir(\P^n_{\k})$
have been actively studied in birational geometry since
the 19th century.
The aim of this work is to contribute to the study of Cremona groups
through a motivic viewpoint 
and bring new methods and results to light.

\ssec{Generating sets of Cremona groups}
\hfill 


We address the problem
about the existence of simple generating sets of Cremona groups.
For the Cremona group of complex projective plane,
it has a well-known generating set due to
M. Noether and Castelnuovo.


\begin{theorem}[Noether--Castelnuovo~{\cite{zbMATH02722071,CastelnuovoCr2}}]
\source{motivic-invariants-birational-maps}
\notready\label{thm:noether}
\source{motivic-invariants-birational-maps}
For any algebraically closed field $\k$,
the Cremona group $\Cr_2(\k)$ is generated by $\PGL_3(\k)$
and the Cremona involution 
\[[X:Y:Z] \mapsto [YZ:XZ:XY].\]
\end{theorem}

When $n \ge 3$, the Cremona groups $\Cr_n(\k)$ 
are much less well understood. 
In terms of the size of generating sets, Hudson and Pan showed that 
for an algebraically closed field $\k$ of characteristic zero,
$\Cr_n(\k)$ is never generated by $\PGL_{n+1}(\k)$
together with any set of transformations of bounded degree 
or countably many elements
when $\k$ is uncountable~\cite{HudsonPan}.
It is therefore an intricate
question in higher dimension whether there exist,
group-theoretically or geometrically,
simple transformations
which generate the entire $\Cr_n(\k)$.
See e.g.~\cite[1.C]{BLZ} for a historical account of this problem,
which has been already mentioned in the 1895 
lectures by Enriques.
Currently, no explicit set of generators
of Cremona groups in dimension $n \ge 3$ is known.





Among the birational automorphisms,
regularizable maps form 
one of the most explicit classes
of elements of $\Cr_n(\k)$.
Recall that an element $\phi \in \Cr_n(\k)$
is called regularizable
if there exists a birational map 
$\alpha: \P^n \dashrightarrow X$
to some variety $X$ 
such that $\alpha \circ \phi \circ \alpha^{-1} \in \Bir(X)$ is a regular automorphism.
See Definition~\ref{Def-reg} for 
a more general definition of 
pseudo-regularizable maps. 
Regular automorphisms of $\P^n$ are regularizable,
and so are elements of $\Aut(U) \subset \Cr_n(k)$
for every open $U \subset \P^n$.
Furthermore, all
elements of $\Cr_n(\k)$ of
finite order are regularizable
by Weil's regularization theorem (see e.g.~\cite[Theorem 1.7]{CheltsovCremona}).
Thus Theorem \ref{thm:noether} implies that $\Cr_2(\C)$
is generated by regularizable elements. 


In higher dimensions, our main result shows that
the contrary holds in most situations and accordingly, 
any set of generators
of $\Cr_n(\k)$
is necessarily quite complicated.

 


\begin{theorem}
\source{motivic-invariants-birational-maps}
\notready\label{thm:main}
\source{motivic-invariants-birational-maps}
In each of
the following cases,
the Cremona group
$\Cr_n(\k)$ is not generated by pseudo-regularizable elements
\textup(which include $\PGL_{n+1}(\k)$
and all 
elements of finite order\textup):
\begin{enumerate}
    \item $n = 3$, and all number fields $\k$, or all function fields $\k$ 
    over a number field, over a finite field or over an algebraically closed field; 
    \item $n \ge 4$, and all fields $\k \subset \C$;
    \item $n \ge 5$, and all infinite fields $\k$.
\end{enumerate}
\end{theorem}



We refer to Theorem \ref{thm:bir-generation} for a more general statement regarding $\Bir(X)$
for classes of varieties $X$ birational to 
$\P^n \times W$ where $W$ is a (separably)
rationally connected variety.
Theorem \ref{thm:main}
disproves a conjecture by Cheltsov in 2004,
which says that 
birational automorphism groups are generated by regularizable
elements \cite[Conjecture 1.12]{CheltsovCremona}.
It also
gives a negative answer to
a more specific
question of Dolgachev \cite[Question on p.1]{Deserti} on whether Cremona groups
are generated by involutions.
We will explain below why certain
types of fields appear in Theorem~\ref{thm:main}.
For now let us note that even though
any field extension $\k \subset \k'$
induces
an embedding $\Cr_n(\k) \subset \Cr_n(\k')$,
there is no straightforward group-theoretic comparison
between $\Cr_n(\k)$ and 
$\Cr_n(\k')$ (in particular $\Cr_n(\ol{\k})$) in
terms of their generating sets.


\ssec{Motivic invariants}
\hfill

Our method of proof relies on a 
motivic 
invariant of birational maps.
Here, by `motivic' we mean
in the sense of
the Grothendieck ring of varieties,
which is one the original
constructions 
of the ring of motives due to Grothendieck, 
and the Burnside ring of Kontsevich--Tschinkel
\cite{KontsevichTschinkel}.
For an $n$-dimensional variety $X$ over $\k$,
this is a group homomorphism
\[
c: \Bir(X) \to \Z[\Bir_{n-1}/\k].
\]
Here $\Bir_{n-1}/\k$ is the set of $(n-1)$-dimensional
birational
equivalence classes of varieties over $\k$.
The invariant $c$ records the birational
classes of 
the exceptional divisors created
and contracted by the birational map.
We prove that in the situations described
in Theorem \ref{thm:main} for $X = \P^n$, 
the invariant 
is non-trivial, 
while it vanishes on all pseudo-regularizable elements.
Furthermore, 
using the explicit non-vanishing of the
invariant $c$ we can 
construct infinitely many
homomorphisms $\Cr_n(\k) \to \Z$ for various
fields $\k$ and $n \ge 3$, see Corollary~\ref{cor:homZ}.
This is in contrast to the 
homomorphisms $\Cr_n(\k) \to \Z/2$ constructed in~\cite{BLZ}.


Our invariant $c$ 
has an enhancement
\[
\ti{c}: \Bir(X) \to \Kt_0(\Var^{\le n-1}/\k),
\]
taking values in the truncated
Grothendieck group of varieties
in such a way that $c$ is the composition
\[
\Bir(X) \overset{\ti{c}}{\to} \Kt_0(\Var^{\le n-1}/\k) \to \Z[\Bir_{n-1}/\k],
\]
where the second homomorphism is the projection to the top dimensional
birational classes. We develop the invariant $\ti{c}$ 
because some proofs are more transparent in this context,
and also because
it provides full control over
the Grothendieck ring of varieties, 
as it determines how
$\Kt_0(\Var^{\le n-1}/\k)$ is mapped to $\Kt_0(\Var^{\le n}/\k)$ 
(see Proposition~\ref{prop:ker-iota}).
None of these constructions require
any form of the
resolution of singularities, and
we work over an arbitrary field
$\k$.

Once the invariants $\ti{c}$ and $c$ are constructed
and their properties settled in Section \ref{sec:invariants}, 
in Section \ref{sec:geometric-constructions}
we introduce and
construct examples of the so-called L-links 
(L stands for the Lefschetz class $\L$
in the Grothendieck ring and L-equivalence \cite{KuznetsovShinder}); 
these links provide a geometric input suitable
for L-equivalence constructions 
and for showing the non-vanishing of $c(\Bir(X))$.
To construct non-trivial L-links, 
we rely on the geometric constructions
of L-equivalence
involving genus one curves
\cite{Mori-Mukai, Crauder-Katz, ShinderZhang}, 
K3 surfaces \cite{HassettLai}, 
and Calabi-Yau threefolds \cite{IMOUG2}.
Even though these constructions already exist in the literature when
$\k = \C$,
much effort goes into generalizing them and checking 
that such L-links $\phi$ are motivically non-trivial,
namely $c(\phi) \ne 0$,
over a more general field $\k$ as in Theorem \ref{thm:main}.

For instance, links involving genus one curves 
become motivically trivial over $\C$, 
but not over
small fields as 
in Theorem \ref{thm:main}(1);
roughly speaking, 
Theorem \ref{thm:main}(1) 
works for these fields
because they have infinitely many
Galois extensions of fixed degree.
To prove Theorem~\ref{thm:main}(2),
we use the construction~\cite{HassettLai},
taking as input a general complex
K3 surface of degree $12$ and
Picard rank one, 
together with some rational points on it.
We extend the construction to $\Q$ nevertheless,
representing K3 surfaces
of degree $12$ 
as hyperplane sections of
a \emph{rational} Fano threefold of degree $12$
defined over $\Q$,
and using Terasoma's argument to prove the existence of K3 surfaces
over $\Q$ of Picard rank one~\cite{Terasoma}
involving Lefschetz pencils.
Finally to prove Theorem \ref{thm:main}(3),
we need to show that the L-links with Calabi-Yau threefolds as centers \cite{IMOUG2}
exist and
can be nontrivial
over any infinite field $\k$.
We have to deal with 
Calabi-Yau varieties 
which are possibly uniruled, and
standard arguments in characteristic zero involving MRC fibrations
do not apply. Instead, we rely on
a weaker statement in the spirit of MRC fibrations
for separably rationally connected varieties,
which we prove in Appendix~\ref{app:MRC}.
We also include
Theorem~\ref{thm:BurtCY},
due to Burt Totaro, showing that
birational
varieties with nef canonical class and 
Picard number
one are isomorphic, in arbitrary
characteristic, which was previously
known in characteristic zero or in positive characteristic in dimension up to three.



\subsection{Relation to other work}

\subsubsection{L-equivalence} 
\hfill

This paper was motivated by 
understanding the geometric meaning
of the information encoded by the Grothendieck ring of varieties and
its variants,
which have played a spectacular role in the recent
results on the specialization of stable rationality \cite{NicaiseShinder} and rationality \cite{KontsevichTschinkel}.
In particular 
we want to understand
L-equivalence \cite{Borisov}, \cite{HassettLai},
\cite{KuznetsovShinder}, \cite{IMOUG2}; our concept of L-links
formalizes the existing geometric constructions leading to L-equivalence.


This is a continuation of our work \cite{LSZ20} where
jointly with Susanna Zimmermann we have studied
the invariant $c(\phi)$ for surfaces $X$ over perfect fields 
using Minimal Model Program.
We showed that
$c(\Bir(X)) = 0$ 
and as a consequence, birational maps between surfaces 
do not produce non-trivial L-equivalence.

On the other hand, 
non-trivial L-equivalence
already exists from dimension $3$ and on,
as we will
explain a construction in Subsection \ref{ssec-dim3} originating from curves in rational threefolds. 
Namely, by Theorem \ref{thm:ell-curves}, we have $\L([C] - [C']) = 0$ for a pair of genus one curves.
This improves the exponent of $\L$ in 
\cite{ShinderZhang} from $\L^4$ to the minimal possible.





\subsubsection{Studying birational maps through the blow up centers} 
\hfill

Making an invariant $c$ from the blow up centers
or the exceptional divisors is similar in spirit 
to the filtration on the Cremona groups by the 
genus of the curves blown up by 
the maps
of threefolds,
considered and studied
by Frumkin \cite{Frumkin-Cremona},
Lamy
\cite{Lamy-Cremona}, Bernardara \cite{Bernardara-Cremona}
and Blanc--Cheltsov--Duncan--Prokhorov \cite{BCDP}.
In fact similar considerations also appear in 
the proof of the Hudson--Pan theorem
\cite{HudsonPan}, which relies on 
the existence of a large
set of possible types of exceptional divisors
for Cremona transformations.
Comparing to the above works,
our approach has the key feature
that the invariant $c$ is a \emph{homomorphism}, which in a way
refines these filtrations.

The relationship between birational maps and the structure
of the Grothendieck ring of varieties also
appears in the work
of Zakharevich in terms of the differentials
in a certain spectral sequence \cite{Zakharevich}
which converges to the higher K-theory groups
lying over
the Grothendieck ring of varieties.
Our invariant $c(\phi)$ coincides with the first
differential from that spectral sequence \cite[Lemma 3.2]{Zakharevich},
and the invariant $\ti{c}(\phi)$ is a certain
enhancement of that map.


On the other hand, relating the Grothendieck
ring of varieties and the Cremona groups is
implicitly suggested in \cite{HassettLai}.
Part of the motivation for our work was to understand
to which extent the blow up centers of a rationality construction $\P^n \dto X$ are determined
by $X$ itself, which was asked in \cite[Introduction]{HassettLai}.
The intrinsic ambiguity in these centers 
leads to the non-vanishing of $c(\Bir(X))$,
which makes our applications to the Cremona groups possible.
Along the way we answer a question
of Hassett and Tschinkel about the 
realization of non-isomorphic genus one curves as
 factorization centers 
for rational threefolds over non-closed fields~\cite[Remark 23 (2)]{HassettTschinkelCycle}, by Theorem \ref{thm:ell-curves} which
was already mentioned above.



\subsubsection{Non-simplicity of the Cremona groups and constructions of non-trivial homomorphisms} 
\hfill

There has been important progress concerning the non-simplicity of the Cremona groups
over the last decade;
see~\cite{Cantat-Cremona} and~\cite[Introduction]{BLZ}
for an excellent overview of the history and recent results.
First of all,
in 2012
Cantat and Lamy proved that 
for any algebraically closed
field $\k$,
$\Cr_2(\k)$ is not a simple group~\cite{MR3037611}.
This result was later extended to $\Cr_2(\k)$ for any field $\k$ by Lonjou~\cite{MR3533276}.
Using the Minimal Model Program and Birkar's boundedness,
Blanc--Lamy--Zimmermann~\cite{BLZ} recently constructed infinitely
many homomorphisms $\Cr_n(\k) \to \Z/2$ ($n \ge 3$) when $\k$ is a subfield of $\C$ (see also~\cite{MR4136435}). 
Their constructions
originate from Sarkisov links and relations between them, and
are of quite different nature to 
the homomorphisms 
$\Cr_n(\k) \to \Z$ we obtain 
in Corollary~\ref{cor:homZ}
through the motivic invariants.

\subsubsection{Generation of Cremona groups
by involutions}
\hfill

As a consequence of Theorem \ref{thm:noether},
$\Cr_2(\C)$ is generated by involutions.
The question whether same is true in higher dimension
was asked by Dolgachev during his series of lectures
in Toulouse in 2016.
Motivated by this question, D\'eserti has proved that $\PGL_{n+1}(\C)$
is generated by involutions \cite[Proposition C]{Deserti}.
Lamy and Schneider have proved
that Cremona groups $\Cr_2(\k)$ over a perfect field $\k$ are generated
by involutions \cite{LamySchneider}, thus providing
a positive answer to Dolgachev's question for $n = 2$.
On the other hand,
Blanc, Schneider and Yasinsky announced 
that birational automorphisms of Severi--Brauer surfaces
are not generated by elements of finite order \cite{BSY},
extending previous results by Shramov \cite{Shramov-SB}.
As a consequence of their result, 
they
are able to give a different proof that $\Cr_n(\C)$
for $n \ge 4$
is not generated by $\PGL_{n+1}(\C)$ and elements of finite order.


\subsection{Notation and conventions}
\label{ss:notation}
\source{motivic-invariants-birational-maps}
\hfill

Unless specified otherwise, we work over an arbitrary field $\k$.
By a variety over $\k$,
we mean
an 
irreducible
and reduced separated
scheme
of finite type over $\k$.
We will use various Grothendieck groups and rings, which we summarize for the convenience of the reader:
\begin{enumerate}
\item[(a)] $\Z[\Bir/\k]$ (resp. $\Z[\Bir_n/\k]$) 
is the free abelian group generated by birational isomorphism classes of varieties 
(resp. varieties of dimension $n$)
over $\k$.
By definition 
$\Z[\Bir/\k] = \bigoplus_{n \ge 0}
\Z[\Bir_n/\k]$. This is the graded Burnside ring of Kontsevich--Tschinkel \cite{KontsevichTschinkel};
\item[(b)] $\Kt_0(\Var/\k)$ 
(resp. $\Kt_0(\Var^{\le n}/\k)$)
is the Grothendieck ring (resp. group)
of varieties (resp. varieties of dimension $\le n$)
over $\k$
modulo cut and paste relations; 
see Subsection~\ref{ssec-tG0} for the precise definitions;
\item[(c)] $\Kt_0(\PPAV/\k)$ is the Grothendieck group of the additive category of 
principally polarized abelian varieties over $\k$
modulo relations $[A \times B] = [A] + [B]$;
\item[(d)] for a profinite group
$G$, $\Kt_0(\Rep(G, \Q_\ell))$
is the Grothendieck group of
the abelian category of
finite-dimensional
continuous $G$-representations 
with coefficients in $\Q_\ell$. See \cite[Section 2.1]{LSZ20} for the details of this construction.
\end{enumerate}
Among these,  (a) and (b)
are central for this paper and
are explained in detail and 
related to each other in Subsection \ref{ssec-tG0}.
The Grothendieck groups (c) and (d)
will be used for some technical purposes
as ``motivic realizations'' of varieties and birational classes,
see in particular \eqref{eq:sigma}, \eqref{eq:j}.
